\chapter*{Annex: Supplementary Evidence and Analytical Material}
%Maynooth complaint formatting 
\addcontentsline{toc}{chapter}{Annex}
\markboth{Annex}{Annex}
%\markboth{}{}

\renewcommand{\thetable}{A.\arabic{table}}
\renewcommand{\thefigure}{A.\arabic{figure}}
\setcounter{table}{0}
\setcounter{figure}{0}


\section*{Thematic Analysis: Methodology and Coding Framework}

The 178 primary sources were subjected to systematic bilateral thematic analysis using two coding approaches. Long-form sources, including documentaries, practitioner interviews and extended combat footage, were segmented into discrete tactical episodes and coded via spreadsheet across source, contextual and analytical fields, including source type, environment, contact type, episode type, troop disposition, survivability, sensing, movement, battle-drill behaviour, enemy and friendly drone presence, cue origin, unit size, initiative, outcome, leadership remarks and thesis implication. Short-form clips were organised directly into thematic categories reflecting the tactical phenomenon depicted: persistent ISR; ISR-corrected fires; battle damage assessment; FPV strikes against armour, vehicles, infantry and shelters; counter-drone action; drone-versus-drone engagement; close-air-support drone employment; UGS-enabled logistics and CASEVAC; armed UGS employment; human factors under drone threat; camouflage and signature management; drone-team survivability and munition preparation; movement under observation; and tactical adaptation under persistent surveillance.

\pagebreak
\section*{Curated Examples of Tactical Behaviour}
The examples below provide a curated evidential register of representative primary-source material.

\begin{itemize}
	\item UAS targeting armoured vehicles: \textcite{TANK_FPV9,TANK_FPV,TANK_FPV3,TANK_FPV4,TANK_FPV6,TANK_FPV9,TANK_FPV8}.
	\item Combined-arms destruction of armoured vehicles cued by and involving drones:  \textcite{TANK_FPV9,TANK_FPV1,TANK_FPV2,TANK_FPV5,TANK_FPV7}.
	\item Flat Ukrainian steppe lending itself to optimal UxS performance: \textcite{ombr72_2023_03_02_funny_fairytale_for_adults} and elucidated by \textcite{JONES_2023}.
	\item Daylight reconnaissance patrol: \textcite{tac_groiup_2026_2_27} in a daylight reconnaissance patrol.
	\item Difficulty for UxS pilots to interpret the terrain through a 2D feed: \textcite{magyar_telegram_FPV_inf_plying_dead}.
	\item Detection being converted into almost immediate lethality: \textcite{azov_media_20240703_no_one_left_to_finish_off}.
	\item Use of UAS for real-time adjusting of fires, battle damage assessment and dynamic tactical decision making: \textcite{azov_media_20240717_three_russian_soldiers_surrender_azov} and \textcite[14:50--20:00]{Wali}.
	\item Point of view experience of dismounted infantry being targeted by an FPVD: \textcite{zandar_2025}.
	\item Being targeted by a CAS: \textcite[1:53--1:58]{Trofimova}.
	\item Vulnerability of entrances to defensive positions to drones and infantry assault: \textcite{225_oshp_2026_03_11_russian_destroy_open,landforces_ukraine_2026_fpv_seen_2026_3_24,tretii_armiiskyi_korpus_2023_3_12_prisoners,azov_media_2025_05_28_assault,mod_russia_telegram_attack_bunker}.
	\item Primary source testimony of defending against FPVDs: \textcite{Tyshchenko_2026_03_20_ukraine_bunker_story}.
	\item Exceptional vulnerability of light forces to UAS: \textcite{azov_media_20240522_enemy_infantry_outside_cover_night_drones,azov_media_20240606_night_elimination_enemy_infantry_kreminna}.
	\item FPV drones increasing infantry lethality at night: \textcite{zandar_2025}.
	\item Waiter FPV drone:  \textcite{rosgvardiya_20250826_kharkiv_steel_brigade_uav_infantry_strike}.
	\item Use of UGS to mitigate human exposure to threats: \textcite{third_army_corps_2025_10_23}.
	\item UGS' greater difficulty than UAS in traversing terrain: \textcite{ombr_63_2026_03_13_casevac}.
	\item Unarmoured CASEVAC UGS variant: \textcite{tretii_armiiskyi_korpus_2023_2_13_CASEVAC}.
	\item UGS as disaggregated support by fire platforms: \textcite{tretii_armiiskyi_korpus_2025_09_02_ground_robot_unit}.
	\item UGS as machine gun platforms during the assault: \textcite{II_mechanized_battalion_ab3_2026_03_22_nrk_turret_forward_edge}.
	\item UGS inducing surrender of the enemy: \textcite{DevDroid_TW762} and \textcite{MODUKR_drone_surrender}.
	\item Russian soldier evading five FPV-OWADs, while being observed by an ISRD: \textcite[9:36--10:38]{fpv_scatter_analysis}. 
	\item Robotics mitigating the vulnerability of light forces to encirclement and bypass: \textcite{Kirichenko_2026B}.
	\item Light infantry assaults remain possible in the most contested areas of the battlefield: \textcite{revanche_telegram_stepnohirsk_2026_2_10}.
	\item A detachment-level assault: \textcite{blue_2_2026_1_13}.
	\item Aggressive and decisive movement through the grey zone: \textcite{Chosen_2024_10_17}.
	\item Commanders perched cybernetically: \textcite[6:33--8:43]{civdiv_6} and \textcite{Dogwoof_2000m_Andriivka}.
	\item UAS used to correct grenade throwing during the battle: \textcite{Chosen_2024_10_17}.
	\item UxS as IEDs: \textcite{civdiv_7,magyar_telegram_iedmaking}, \textcite{civdiv}, \textcite[7:30--7:37]{civdiv_2} and discussed with \textcite{danny_participant_5}.
	\item Defence of a moving vehicle from an FPV-OWAD by a shotgun: \textcite[0:16--0:50]{civdiv_11} illustrates.
	\item Live reflection on one's own battle via UAS footage: \parencite{Dave_2024_3_22,hallatt_2024_03_13_1}.
	\item Short, accessible lessons-learned material: \textcite{uKR_LESSONS_LEARNED_1}.
	\item Thermal camouflage held above the body in a manner analogous to the scrim umbrella: \index{Camouflage!Scrim umbrella} \textcite[0:00--0:35]{civdiv_11}.
	\item Vulnerability of defensive firing ports to FPV-OWADs: \textcite{azov_media_2025_05_28_assault,mod_russia_telegram_attack_bunker}.
\end{itemize}

 
 

\pagebreak
\section*{Translations of Non-English Source Author Names}

These tables provide the original non-English source author names and the English forms used for in-text citation and bibliography entries throughout the thesis.
\linebreak

\begin{table}[H]
	\centering
	\small
	\begin{tabular}{p{0.44\textwidth} p{0.44\textwidth}}
		\toprule
		Original Author Name & English Author Name Used in Thesis \\
		\midrule
		\multicolumn{2}{c}{\textbf{Russian sources}} \\
		\midrule
		83 Уссурийская & 83rd Ussuriysk Air Assault Brigade \\
		Министерство обороны Российской Федерации & Russian Ministry of Defence \\
		Родная 98-я ВДД & 98th Guards Airborne Division \\
		Росгвардия & National Guard of Russia \\
		\bottomrule
	\end{tabular}
	\caption[Translations of Russian Source Author Names]{Translations of Russian Source Author Names Used in the Thesis}
	\label{tab:russian-source-name-translations}
\end{table}

\begin{table}[H]
	\centering
	\small
	\begin{tabular}{p{0.44\textwidth} p{0.44\textwidth}}
		\toprule
		Original Author Name & English Author Name Used in Thesis \\
		\midrule
		\multicolumn{2}{c}{\textbf{Ukrainian sources}} \\
		\midrule
		2-й механізований батальйон 3 ОШБр & 2nd Mechanised Battalion, 3rd Assault Brigade \\
		225 ОШП & 225th Separate Assault Regiment \\
		47 окрема механізована бригада ``Маґура'' & 47th Separate Mechanised Brigade ``Magura'' \\
		63 бригада \textbar{} Сталеві Леви & 63rd Brigade \textbar{} Steel Lions \\
		66 ОМБр ім.\ князя Мстислава Хороброго & 66th Separate Mechanised Brigade named after Prince Mstyslav the Brave \\
		72 ОМБр ім.\ Чорних Запорожців & 72nd Separate Mechanised Brigade named after the Black Zaporozhians \\
		Азов Media \textbar{} Azov Media & Azov Media \\
		Військово-Морські Сили ЗС України \textbar{} UA Navy & Ukrainian Navy \\
		Держприкордонслужба & State Border Guard Service of Ukraine \\
		Kraken \textbar{} ПБС 3 армійського корпусу & Kraken \textbar{} Special Purpose Battalion, 3rd Army Corps \\
		Мілітарний & Militarnyi \\
		Об'єднана штурмова бригада НПУ ``Лють'' & Lyut Joint Assault Brigade, National Police of Ukraine \\
		Птахи люті & Birds of Fury \\
		Служба безпеки України & Security Service of Ukraine \\
		Сухопутні війська \textbar{} UA Land Forces & Ukrainian Land Forces \\
		Тактична група ``Реванш'' & Revanche Tactical Group \\
		Третій армійський корпус & Third Army Corps \\
		\bottomrule
	\end{tabular}
	\caption[Translations of Ukrainian Source Author Names]{Translations of Ukrainian Source Author Names Used in the Thesis}
	\label{tab:ukrainian-source-name-translations}
\end{table}

\pagebreak

%consider a graphical depiction --- pie chart per molloy_2024
\pagebreak
\small
\renewcommand{\arraystretch}{1.1}
\setlength{\tabcolsep}{4pt}

\pagebreak
\section*{Primary Source Interviews}

Primary source interviews conducted for this thesis are anonymised as Participants 01--10. Publicly available primary-source interviews are cited using the name, callsign, handle or alias under which the source is publicly identified.

\begin{tabularx}{\textwidth}{@{} l l X @{}}
	\toprule
	\textbf{Participant} & \textbf{Background} & \textbf{Date} \\
	\midrule
	P01 & Russo--Ukrainian War veteran, private & 23 Jan 2026 \\
	P02 & War correspondent; analyst & 28 Jan 2026 \\
	P03 & Irish military general officer & 5 Feb 2026 \\
	P04 & Russo--Ukrainian War veteran, junior NCO & 6 Feb 2026 \\
	P05 & Senior EU civil servant & 6 Feb 2026 \\
	P06 & Russo--Ukrainian War veteran, private & 6 Feb 2026 \\
	P07 & USSOCOM officer & 23 Feb 2026 \\
	P08 & Defence researcher & 12 Mar 2026 \\
	P09 & Russo--Ukrainian War veteran, junior NCO & 4 Apr 2026 \\ 
	P10 & Defence academic & 4 Apr 2026  \\
	%P13 & Irish military general officer & Planned --- likely \\
	\bottomrule
\end{tabularx}
